\chapter{Open Source Foundation and Licensing}
\label{ch:foundation}

\section{The commercial constraint}

Every dependency in Xyra must satisfy one rule: it must permit distribution of a binary
whose corresponding source is not disclosed. In practice this means MIT, Apache~2.0, BSD,
ISC, Zlib, Unlicense or CC0. It excludes GPL and LGPL in linked form, and it excludes MPL~2.0
as well, for the reason given below.

This rule is what separates the standalone product from the current fork. It is not a
philosophical preference, it is the precondition for a commercial edition.

The allow list is exactly MIT, Apache-2.0, BSD-2-Clause, BSD-3-Clause, ISC, Zlib, Unlicense
and CC0-1.0. MPL-2.0 is \req{not} on it. Its copyleft is file scoped rather than project
scoped, which makes it survivable in principle, but a per file disclosure obligation inside a
commercial binary is a compliance burden with no compensating benefit when a permissive
alternative exists. A dependency outside the list enters only through an explicit entry in the
exception list inside the license check script, which puts the decision in front of a reviewer
rather than inside a transitive dependency nobody read.

\section{Verified license position}

The table below was compiled from the crates.io registry metadata and from the license
declarations inside each project, rather than from documentation or memory. Every entry was
checked while writing this document.

\begin{longtable}{@{}p{0.24\textwidth} p{0.30\textwidth} p{0.34\textwidth}@{}}
\toprule
\textbf{Component} & \textbf{License} & \textbf{Role in Xyra} \\
\midrule
\endfirsthead
\toprule
\textbf{Component} & \textbf{License} & \textbf{Role in Xyra} \\
\midrule
\endhead
\multicolumn{3}{@{}l}{\textit{Interface and rendering}}\\
gpui & Apache-2.0 & GPU accelerated UI framework, element tree, layout, input \\
wgpu & MIT OR Apache-2.0 & Graphics abstraction, fallback rendering path \\
vello & Apache-2.0 OR MIT & Vector rendering, alternative compositor path \\
cosmic-text & MIT OR Apache-2.0 & Text shaping and layout \\
swash & Apache-2.0 OR MIT & Font parsing, glyph rasterisation \\
\midrule
\multicolumn{3}{@{}l}{\textit{Editor core}}\\
ropey & MIT OR Apache-2.0 & Rope data structure for the text buffer \\
tree-sitter & MIT & Incremental parsing, concrete syntax trees \\
tree-sitter-highlight & MIT & Syntax highlighting over the parse tree \\
similar & Apache-2.0 & Text diffing for review surfaces \\
imara-diff & Apache-2.0 & High performance line diffing for large files \\
\midrule
\multicolumn{3}{@{}l}{\textit{Language intelligence}}\\
lsp-types & MIT & Language server protocol type definitions \\
tower-lsp or hand rolled client & MIT & Language server transport and lifecycle \\
\midrule
\multicolumn{3}{@{}l}{\textit{Project and system}}\\
gix (gitoxide) & MIT OR Apache-2.0 & Pure Rust git implementation \\
ignore & Unlicense OR MIT & Gitignore aware directory traversal \\
globset & Unlicense OR MIT & Glob matching for path rules \\
notify & CC0-1.0 & Filesystem watching \\
alacritty\_terminal & Apache-2.0 & Terminal emulation backend \\
tokio & MIT & Asynchronous runtime for the service layer \\
dashmap & MIT & Concurrent maps for shared indices \\
\midrule
\multicolumn{3}{@{}l}{\textit{Retrieval and storage}}\\
tantivy & MIT & Full text and structural search index \\
usearch & Apache-2.0 & Approximate nearest neighbour vector index \\
fastembed & Apache-2.0 & Local embedding generation \\
candle-core & MIT OR Apache-2.0 & Local model execution for embeddings and rerankers \\
tiktoken-rs & MIT & Token counting for budget accounting \\
redb & MIT OR Apache-2.0 & Embedded key value store for the durable state \\
rusqlite & MIT & Relational store for the graph and journal \\
zstd & MIT & Compression of context artefacts and caches \\
\midrule
\multicolumn{3}{@{}l}{\textit{Agent integration}}\\
agent-client-protocol & Apache-2.0 & Protocol for external agent processes \\
rmcp & Apache-2.0 & Model context protocol client and server \\
wasmtime & Apache-2.0 WITH LLVM-exception & Sandboxed extension execution \\
\midrule
\multicolumn{3}{@{}l}{\textit{Interface support}}\\
fuzzy-matcher & MIT & Fuzzy matching for pickers and the palette \\
\bottomrule
\caption{Verified license position of the proposed dependency set.}
\end{longtable}

\section{What this table means}

Two conclusions follow directly.

First, \req{a complete IDE can be assembled without touching a single copyleft dependency}.
The UI framework, the rope, the parser, the terminal, the git implementation, the language
server client, the vector index and the extension sandbox are all permissive. There is no
component in the critical path that forces source disclosure.

Second, \req{the highest leverage reuse is gpui}. It is Apache~2.0, it is production proven
in a shipping editor, and it solves the single hardest problem in this project, which is a
GPU accelerated UI framework with correct text rendering and input handling on three
platforms. Writing that from scratch would consume the majority of the engineering budget
and produce a worse result. Using it is not a compromise on independence, because Apache~2.0
imposes no disclosure obligation and no upstream coupling beyond the version we pin.

\section{Dependency strategy}

\begin{description}[style=nextline,leftmargin=0pt]
  \item[Vendor and pin the UI framework.] gpui is vendored into the repository at a pinned
  revision under its Apache~2.0 license with attribution preserved. Upstream changes are
  merged deliberately, never automatically. This removes the rebase treadmill that
  characterises the fork.

  \item[Wrap every third party boundary.] Editor code never calls gpui directly. It calls
  \ctl{xyra-ui}, a thin internal crate that re-exports the element model we actually use.
  The same applies to the parser, the git layer and the terminal. The cost is one indirection
  and the benefit is that any dependency can be replaced without touching product code.

  \item[Prefer boring dependencies.] Fewer, larger, well maintained crates over many small
  ones. Every dependency is a license question, a supply chain question and a maintenance
  question.

  \item[No network dependency at build time.] The build must succeed from a vendored source
  tree with no registry access, which is a hard requirement for reproducibility and for
  enterprise customers with air gapped environments.
\end{description}

\section{License hygiene in the repository}

\begin{itemize}
  \item A machine readable manifest at \ctl{licenses/manifest.toml} lists every dependency,
  its license, its source and the reason it is present.
  \item A continuous integration job fails the build when a dependency appears whose license
  is not on the allow list, using cargo-deny or an equivalent check.
  \item Vendored third party source lives under \ctl{third\_party/} with its original license
  file intact and unmodified.
  \item The notice file shipped with the product aggregates every attribution required by
  the Apache~2.0 and MIT terms.
\end{itemize}

\section{Reference implementations worth studying}

These projects are not dependencies. They are prior art whose architecture repays study
before writing the equivalent subsystem. All are permissively licensed, so reading them
carries no contamination risk, and porting an approach from them is legally clean when the
license is respected and attribution given.

\begin{table}[h]
\centering
\small
\begin{tabularx}{\textwidth}{l l X}
\toprule
\textbf{Project} & \textbf{License} & \textbf{What to learn from it} \\
\midrule
Helix & MPL-2.0 & Multi cursor and selection model, tree-sitter integration depth \\
Lapce & Apache-2.0 & Rope plus WASM plugin boundary, multi window process model \\
Zed gpui & Apache-2.0 & Element tree, layout engine, GPU text pipeline \\
Alacritty & Apache-2.0 & Terminal grid, escape sequence handling, performance discipline \\
Tantivy & MIT & Index segment design, query planning \\
Wasmtime & Apache-2.0 & Capability based sandboxing for untrusted extensions \\
Ripgrep and ignore & MIT & Directory traversal that respects ignore semantics at speed \\
\bottomrule
\end{tabularx}
\caption{Permissively licensed prior art.}
\end{table}
